% ------------------------------------------------------------------------------------------------
% 4.1. Emparelhamento de cardinalidade máxima
% ------------------------------------------------------------------------------------------------
\subsection{Emparelhamento de cardinalidade máxima}
\label{sec:emp_card_max_bip}

\subsubsection{Descrição do problema}

Dado um grafo não-direcionado $G=(V,E)$, o problema de emparelhamento de cardinalidade máxima (MCM) busca encontrar um emparelhamento $M \subseteq E$ tal que o número de arestas em $M$ seja maximizado, ou seja, $|M|$ é o maior entre todos os emparelhamentos possíveis em $G$.

Um caso especial, mas importante, do MCM é o \textbf{Emparelhamento perfeito}, que é um emparelhamento onde todos os vértices do grafo são cobertos por exatamente uma aresta do emparelhamento. Para que um emparelhamento perfeito exista, o grafo deve ter um número par de vértices e o emparelhamento máximo deve ter cardinalidade igual a $|V|/2$ \cite{jungnickel}.

\subsubsection{Propriedades}

Antes de continuar para os algoritmos e soluções, é necessário entender algumas propriedades importantes

\paragraph{Emparelhamento máximo vs maximal}
\label{p:emp_max_maximal}

Um emparelhamento máximo é aquele que contém o maior número possível de arestas, enquanto um emparelhamento maximal é aquele que não pode ser aumentado adicionando mais arestas, mas não necessariamente é o maior possível \cite{manber}.

\paragraph{Caminhos aumentantes}
\label{p:caminhos_aumentantes}

Um caminho aumentante $P$ em um grafo $G$ é chamado de caminho aumentante se $P$ tem tamanho ímpar, começa e termina em vértices livres (não cobertos por nenhuma aresta do emparelhamento atual) e alterna entre arestas que pertencem a $M$ e arestas que não pertencem a $M$. A existência de um caminho aumentante indica que o emparelhamento atual não é máximo \cite{schrijver}

\paragraph{Teorema de Kőnig}
\label{p:teorema_konig}

O teorema de Kőnig estabelece uma relação fundamental entre emparelhamentos e coberturas em grafos bipartidos. Ele afirma que, em qualquer grafo bipartido, o tamanho do maior emparelhamento é igual ao tamanho da menor cobertura de vértices \cite{jungnickel}.

A cobertura de vértices é um conjunto de vértices tal que cada aresta do grafo é incidente a pelo menos um vértice desse conjunto. O teorema de Kőnig é uma ferramenta poderosa para resolver problemas de emparelhamento, utilizado nos algoritmos de Hopcroft-Karp e no método Húngaro. A seguir será detalhado um método algorítmico para encontrar a cobertura de vértices mínima a partir de um emparelhamento máximo.

O método é definido pelos seguintes passos \cite{jungnickel}:
\begin{enumerate}
    \item Identifica-se o conjunto de vértices livres (não saturados por $M$) na partição esquerda $U$.
    \item A partir desses vértices, inicia-se uma travessia no grafo (utilizando Busca em Largura ou Profundidade) com regras estritas de alternância: transições do conjunto $U$ para $V$ devem ocorrer exclusivamente através de arestas que \textbf{não} pertencem ao emparelhamento ($E \setminus M$), enquanto as transições de $V$ de volta para $U$ devem utilizar obrigatoriamente as arestas do emparelhamento ($M$).
    \item Registra-se em um conjunto $Z$ todos os vértices que foram alcançados durante essa busca, tanto da partição $U$ quanto da partição $V$.
    \item A Cobertura Mínima de Vértices $C$ é construída selecionando-se os vértices de $U$ que \textbf{não} foram visitados e os vértices de $V$ que \textbf{foram} visitados. Matematicamente, a cobertura é dada por $C = (U \setminus Z) \cup (V \cap Z)$.
\end{enumerate}

Essa extração algorítmica possui complexidade linear em relação ao tamanho do grafo, $O(V + E)$, visto que requer apenas uma travessia simples após o cálculo do emparelhamento. Sua implementação é uma sub-rotina fundamental para métodos exatos de problemas mais complexos, como a atualização de limiares no Método Dual para o problema de atribuição gargalo \cite{cormen}.

\paragraph{Teorema de Hall}

O Teorema de Hall, frequentemente referido como Teorema do Casamento de Hall, é outro resultado estrutural para grafos bipartidos. Enquanto o teorema de König determina a cardinalidade do emparelhamento máximo, o teorema de Hall estabelece a condição necessária e suficiente para a existência de um emparelhamento que cubra integralmente um dos subconjuntos da partição \cite{schrijver}. 

Para um grafo bipartido $G=(U \cup V, E)$, o teorema afirma que existe um emparelhamento cobrindo todo o conjunto $U$ se, e somente se, para qualquer subconjunto $S \subseteq U$, o tamanho de sua vizinhança $N(S)$ (o conjunto de vértices em $V$ adjacentes a pelo menos um vértice em $S$) for maior ou igual ao tamanho de $S$. Matematicamente, exige-se que $|N(S)| \ge |S|$ para todo $S \subseteq U$ \cite{jungnickel}.  Essa propriedade, conhecida como Condição de Hall, é essencial para provar matematicamente a viabilidade de atribuições e identificar gargalos estruturais que impedem a formação de um emparelhamento perfeito, relacionando-se intimamente com os conceitos de cortes mínimos em redes de fluxo.

\paragraph{Teorema de Gallai e Conjunto Independente Máximo}

Um conjunto independente em um grafo é um subconjunto de vértices no qual não existem duas extremidades adjacentes, ou seja, nenhuma aresta conecta dois vértices desse conjunto. O problema de encontrar o Conjunto Independente Máximo (\textit{Maximum Independent Set} - MIS) é classicamente conhecido por pertencer à classe NP-Difícil para grafos gerais \cite{garey-jhonson}. No entanto, quando o grafo é restrito à classe dos grafos bipartidos, o problema pode ser resolvido em tempo polinomial graças à sua profunda relação com os emparelhamentos e coberturas.

A base dessa resolução é o Teorema de Gallai, proposto em 1959. O teorema estabelece que, para qualquer grafo $G=(V, E)$ que não possua vértices isolados, a soma do tamanho da Cobertura Mínima de Vértices (MVC) com o tamanho do Conjunto Independente Máximo (MIS) é sempre igual ao número total de vértices do grafo \cite{schrijver}. 

Matematicamente:

\begin{equation}
    |MVC| + |MIS| = |V|
\end{equation}

Pelo Teorema de König, detalhado anteriormente, sabemos que em grafos bipartidos o tamanho da Cobertura Mínima de Vértices é exatamente igual ao tamanho do Emparelhamento Máximo ($|MVC| = |M|$). Consequentemente, o tamanho do Conjunto Independente Máximo pode ser deduzido diretamente do emparelhamento: $|MIS| = |V| - |M|$ \cite{jungnickel}.

Além da equivalência de cardinalidade, o Teorema de Gallai também fornece a dualidade estrutural: o Conjunto Independente Máximo é exatamente o complemento da Cobertura Mínima de Vértices. Portanto, para construir algoritmicamente o Conjunto Independente Máximo em um grafo bipartido, basta executar o seguinte procedimento:

\begin{enumerate}
    \item Calcula-se o Emparelhamento Máximo utilizando algoritmos como o de Hopcroft-Karp.\
    \item Extrai-se a Cobertura Mínima de Vértices $C$ através do método de busca por caminhos alternantes a partir dos vértices livres, definido na seção \ref{p:teorema_konig}.\
    \item O Conjunto Independente Máximo $I$ será composto por todos os vértices do grafo que não pertencem à cobertura. Ou seja, $I = V \setminus C$.
\end{enumerate}

\subsubsection{Algoritmos exatos}
\label{sec:algos_exatos_emp_card_max}

Por ser em um ambiente mais restrito, o problema de emparelhamento de cardinalidade máxima tem diversas soluções em tempo polinomial. A seguir listadas algumas delas:

\paragraph{Algoritmo do caminho aumentante}

O algoritmo de caminho aumentante se baseia em um lema proposto por Claude Berge, que afirma que um grafo G é máximo se e somente se não existe nenhum caminho aumentante em G. Este algoritmo é uma implementação direta deste lema, onde se busca repetidamente por caminhos aumentantes e se aumenta o emparelhamento até que nenhum caminho aumentante possa ser encontrado \cite{lawler}.

Começamos com um emparelhamento vazio $M=\emptyset$. Escolhemos um vértice livre $u$ no lado esquerdo do grafo bipartido e tentamos encontrar um caminho aumentante $P$ que começa em $u$ e termina em um vértice livre no lado direito do grafo, respeitando a regra da alternância entre arestas em $M$ e arestas fora de $M$. Se um caminho aumentante for encontrado, atualizamos o emparelhamento $M$ invertendo as arestas ao longo do caminho $P$. Repetimos esse processo até que não seja possível encontrar mais caminhos aumentantes. O emparelhamento resultante será o emparelhamento de cardinalidade máxima \cite{halim}.

Em seguida o algoritmo em pseudo-código:

\begin{algorithm}[H]
    \DontPrintSemicolon
    \SetAlgoLined

    % Variáveis globais
    \textbf{Global:} $match$ (vetor de emparelhamento), $vis$ (vetor de visitados), $AdjList$\;
    \bigskip

    % Procedimento Augment (DFS)
    \Fn{Aug($l$)}{
        \If{$vis[l] = 1$}{
            \KwRet 0\;
        }
        $vis[l] \leftarrow 1$\;
        
        \For{each $r \in AdjList[l]$}{
            % A lógica do caminho aumentante
            \If{$match[r] = -1$ \KwOr Aug($match[r]$)}{
                $match[r] \leftarrow l$\;
                \KwRet 1 \textbf{[comment: found 1 matching]}\;
            }
        }
        \KwRet 0 \textbf{[comment: no matching]}\;
    }

    \bigskip

    % Procedimento Principal
    \Fn{MCBM\_Algorithm($n$)}{
        $MCBM \leftarrow 0$\;
        $match \leftarrow -1$ (initialize for all $V$)\;
        
        \For{$l \leftarrow 0$ \KwTo $n-1$}{
            $vis \leftarrow 0$ (reset before each recursion)\;
            $MCBM \leftarrow MCBM + \text{Aug}(l)$\;
        }
        
        \KwRet $MCBM$\;
    }
    
    \caption{Algoritmo de caminhos aumentantes (MCBM) \cite{halim}}
    \label{alg:mcbm_augmenting}
\end{algorithm}

Como o algoritmo repete a busca por caminhos aumentantes para cada vértice do lado esquerdo do grafo, o tempo de execução total do algoritmo é $O(VE)$, onde $V$ é o número de vértices e $E$ é o número de arestas no grafo bipartido \cite{halim}.

\paragraph{Redução ao problema de fluxo máximo}

O problema de emparelhamento de cardinalidade máxima em grafos bipartidos pode ser eficientemente resolvido através de uma redução ao problema de fluxo máximo em redes. A ideia central é construir uma rede de fluxo a partir do grafo bipartido original, onde cada aresta do grafo bipartido é convertida em uma aresta com capacidade unitária na rede de fluxo. \cite{lawler}

Primeiro, adicionamos um vértice fonte $s$ e um vértice sumidouro $t$ à rede. Conectamos o vértice fonte $s$ a todos os vértices do conjunto esquerdo do grafo bipartido com arestas de capacidade 1. Em seguida, conectamos todos os vértices do conjunto direito do grafo bipartido ao vértice sumidouro $t$, também com arestas de capacidade 1. As arestas entre os conjuntos esquerdo e direito do grafo bipartido são mantidas na rede de fluxo, cada uma com capacidade 1 \cite{halim}.

Com isso, temos agora um grafo de fluxo onde o objetivo é encontrar o fluxo máximo do vértice fonte $s$ para o vértice sumidouro $t$. O valor do fluxo máximo encontrado nesta rede corresponde ao tamanho do emparelhamento máximo no grafo bipartido original. \cite{lawler}. Para encontrar o fluxo máximo, podemos utilizar algoritmos clássicos como o de Edmonds-Karp.

\paragraph{Algoritmo de Edmonds-Karp}
\label{p:edmonds-karp}

O algoritmo de \textbf{Edmonds-Karp} é o principal método para resolver esse tipo de problema. Ele é uma especialização do método de Ford-Fulkerson descrito na seção \ref{p:ford_fulkerson} que utiliza uma busca em largura (BFS) para encontrar sempre o caminho aumentante mais curto em número de arestas. Tem complexidade de tempo $O(VE^2)$:

\begin{algorithm}[H]
    \DontPrintSemicolon
    \SetAlgoLined

    \textbf{Global:} $res$ (matriz de capacidade residual), $mf, f, s, t$, $p$ (vetor de pais da BFS)\;
    \bigskip

    \Fn{BFS($s, t$)}{
        $dist \leftarrow \infty; \quad dist[s] \leftarrow 0$\;
        $Q \leftarrow \{s\}; \quad p \leftarrow \text{nil}$ (init vector with -1)\;
        \While{$Q \neq \emptyset$}{
            $u \leftarrow Q.\text{front}()$; $Q.\text{pop}()$\;
            \If{$u = t$}{
                \KwBreak (stop BFS)\;
            }
            \For{each $v \in V$}{
                \If{$res_{u, v} > 0$ \KwAnd $dist[v] = \infty$}{
                    $dist[v] \leftarrow dist[u] + 1$\;
                    $Q.\text{push}(v)$\;
                    $p[v] \leftarrow u$\;
                }
            }
        }
    }

    \bigskip

    \Fn{Augment($v, \text{minEdge}$)}{
        \eIf{$v = s$}{
            $f \leftarrow \text{minEdge}$\;
            \KwRet\;
        }{
            \If{$p[v] \neq -1$}{
                Augment($p[v], \min(\text{minEdge}, res_{p[v], v})$)\;
                $res_{p[v], v} \leftarrow res_{p[v], v} - f$\;
                $res_{v, p[v]} \leftarrow res_{v, p[v]} + f$\;
            }
        }
    }

    \bigskip

    \Fn{EdmondsKarp($s, t$)}{
        $mf \leftarrow 0$\;
        \While{\KwTrue}{
            $f \leftarrow 0$\;
            BFS($s, t$); \AlgComment{fills the parent vector $p$}\;
            Augment($t, \infty$); \AlgComment{find min edge weight $f$ in path}\;
            \If{$f = 0$}{
                \KwBreak (terminate)\;
            }
            $mf \leftarrow mf + f$\;
        }
        \KwRet $mf$\;
    }

    \caption{Algoritmo de Edmonds-Karp \cite{halim}}
    \label{alg:edmonds_karp}
\end{algorithm}

O pseudocódigo divide-se em três procedimentos:

\begin{itemize}
    \item \textbf{BFS:} Busca em largura a partir de $s$ na rede residual, percorrendo apenas arestas com capacidade residual positiva. Ao alcançar cada vértice pela primeira vez, registra seu pai em $p$; como a busca é em largura, o caminho registrado até $t$ é o de menor número de arestas.
    \item \textbf{Augment:} Percorre o caminho encontrado de $t$ até $s$, seguindo o vetor de pais. Na volta da recursão, $f$ guarda a menor capacidade residual do caminho (o gargalo), que é subtraída das arestas no sentido do caminho e somada às arestas reversas.
    \item \textbf{EdmondsKarp:} Repete a busca e o aumento enquanto existir caminho de $s$ a $t$. Quando a BFS não alcança $t$, nenhum fluxo é enviado ($f = 0$) e o fluxo acumulado $mf$ é máximo.
\end{itemize}

\paragraph{Algoritmo de Dinic}

Outra opção para resolver o problema de fluxo máximo é o algoritmo de Dinic, que é mais eficiente em muitos casos práticos, especialmente em grafos densos. O algoritmo de Dinic utiliza uma combinação de buscas em largura (BFS) para construir níveis na rede residual e buscas em profundidade (DFS) para encontrar caminhos aumentantes dentro desses níveis. O tempo de execução do algoritmo de Dinic é $O(V^2 E)$ em redes gerais e cai para $O(E \sqrt{V})$ em redes de capacidade unitária, como a obtida na redução do emparelhamento bipartido, o que o torna uma escolha preferida para muitos problemas de fluxo máximo \cite{halim}.

Este trabalho não tem uma implementação do algoritmo de Dinic. Ao invés disso, será apresentado o algoritmo de Hopcroft-Karp, que é uma variação especializada do algoritmo de Dinic para o problema de emparelhamento em grafos bipartidos, oferecendo uma solução mais eficiente para este caso específico.

\paragraph{Algoritmo de Hopcroft-Karp}

O algoritmo de Hopcroft-Karp é um algoritmo eficiente para encontrar o emparelhamento máximo em grafos bipartidos. Ele melhora o desempenho do algoritmo de caminhos aumentantes ao encontrar múltiplos caminhos aumentantes em cada iteração, em vez de apenas um. O tempo de execução do algoritmo de Hopcroft-Karp é $O(E \sqrt{V})$, tornando-o significativamente mais rápido do que o algoritmo de caminhos aumentantes simples, especialmente em grafos grandes \cite{halim}.

O algoritmo consiste em duas fases principais: a fase de construção de níveis e a fase de busca de caminhos aumentantes. Na fase de construção de níveis, uma busca em largura (BFS) é realizada a partir dos vértices livres no lado esquerdo do grafo bipartido para construir uma camada de níveis que ajuda a identificar os caminhos aumentantes mais curtos. Na fase de busca de caminhos aumentantes, uma busca em profundidade (DFS) é realizada para encontrar todos os caminhos aumentantes possíveis dentro da camada de níveis construída na fase anterior. Esses caminhos são então usados para aumentar o emparelhamento \cite{halim}.

A seguir, o pseudo-código do algoritmo de Hopcroft-Karp:

\begin{algorithm}[H]
    \DontPrintSemicolon
    \SetAlgoLined

    \caption{Algoritmo de Hopcroft-Karp \cite{burkard}}
    \label{alg:hopcroft_karp}

    % MAIN ALGORITHM
    \KwLet $G = (U, V; E)$ be a bipartite graph with initial (possibly empty) matching $M$;\;
    \KwLet $U_0$ contain all unmatched vertices of $U$;\label{hk:u0}\;
    \KwLet $V_0$ contain all unmatched vertices of $V$;\label{hk:v0}\;
    
    \While{$U_0 \neq \emptyset$\label{hk:phase}}{
        Construct\_layered\_graph;\label{hk:callbuild}\;
        Find\_set\_$\Delta M$;\label{hk:callfind}\;
        $M := M \ominus \Delta M$;\label{hk:symdiff}\;
        update sets $U_0, V_0$ of unmatched vertices\label{hk:upd}\;
    }

    \bigskip

    % PROCEDURE 1: CONSTRUCT LAYERED GRAPH
    \Fn{Construct\_layered\_graph}{
        $L_0 := U_0, k^* := k := 0$;\label{hk:l0}\;
        \While{$L_k \neq \emptyset$}{
            \For{each $i \in L_k$}{
                $N(i) := \{ j : [i, j] \in E \setminus M, j \notin L_1 \cup L_3 \cup \dots \cup L_{k-1} \}$;\label{hk:ni}
            }
            $L_{k+1} := \bigcup_{i \in L_k} N(i)$;\;
            
            \If{$L_{k+1} = \emptyset$\label{hk:stop}}{
                \KwStop \AlgComment{$M$ is a maximum matching}
            }
            
            \eIf{$L_{k+1} \cap V_0 \neq \emptyset$\label{hk:found}}{
                $k^* := k + 1$;\;
                $L_{k+2} := \emptyset$\;
            }{
                $L_{k+2} := \{ \overline{i} : [\overline{i}, j] \in M, j \in L_{k+1} \}$;\label{hk:lk2}
            }
            $k := k + 2$\;
        }
    }

    \bigskip

    % PROCEDURE 2: FIND SET DELTA M
    \Fn{Find\_set\_$\Delta M$}{
        \AlgComment{find a maximal set $\Delta M$ of disjoint augmenting paths}\;
        delete all vertices in $L_{k^*} \setminus V_0$;\label{hk:del}\;
        mark all remaining vertices as unscanned;\;
        $k := 1$; \AlgComment{path counter}\;
        
        \While{$L_0 \neq \emptyset$}{
            choose $x_0 := i \in L_0$ and delete it in $L_0$;\label{hk:x0}\;
            $\ell := 0$;\;
            
            \While{$\ell \ge 0$}{
                \While{$x_\ell$ has an unscanned neighbor in $L_{\ell+1}$\label{hk:adv}}{
                    choose unscanned neighbor $x_{\ell+1}$;\;
                    mark $x_{\ell+1}$ as scanned;\;
                    $\ell := \ell + 1$;\;
                    
                    \If{$\ell = k^*$\label{hk:path}}{
                        $P_k := (x_0, x_1, \dots, x_{k^*})$;\;
                        $k := k + 1$\;
                    }
                }
                \eIf{$\ell < k^*$\label{hk:back}}{
                    $\ell := \ell - 1$\;
                }{
                    $\ell := -1$\;
                }
            }
        }
        \KwRet $\Delta M := (P_1, P_2, \dots, P_{k-1})$\;
    }

\end{algorithm}

O pseudocódigo é organizado em um laço principal e dois procedimentos, explicados a seguir com referência às linhas do Algoritmo~\ref{alg:hopcroft_karp}:

\begin{itemize}
    \item \textbf{Laço principal (linhas~\ref{hk:u0}--\ref{hk:upd}):} $U_0$ e $V_0$ guardam os vértices livres de cada lado (linhas~\ref{hk:u0} e~\ref{hk:v0}). Cada iteração do laço da linha~\ref{hk:phase} é uma \textit{fase}: constrói-se o grafo em camadas (linha~\ref{hk:callbuild}), encontra-se um conjunto maximal $\Delta M$ de caminhos aumentantes mínimos e disjuntos (linha~\ref{hk:callfind}) e o emparelhamento é atualizado pela diferença simétrica $M \ominus \Delta M$ (linha~\ref{hk:symdiff}). Ao longo de cada caminho, as arestas emparelhadas e as não emparelhadas trocam de papel, de modo que $|M|$ aumenta em $|\Delta M|$.
    \item \textbf{Construção do grafo em camadas:} É uma busca em largura que parte de todos os vértices livres de $U$ ao mesmo tempo, que formam a camada $L_0$ (linha~\ref{hk:l0}). As camadas de índice ímpar contêm vértices de $V$, alcançados por arestas \textit{não} emparelhadas (linha~\ref{hk:ni}); a restrição $j \notin L_1 \cup \dots \cup L_{k-1}$ impede que um vértice seja revisitado. Se nenhum vértice novo é alcançado (linha~\ref{hk:stop}), não há caminho aumentante e, pelo lema de Berge, $M$ é máximo. Se a nova camada contém um vértice livre de $V$ (linha~\ref{hk:found}), foi encontrado o comprimento $k^*$ dos caminhos aumentantes mais curtos, e a construção termina. Caso contrário, a camada seguinte é formada pelos parceiros, em $M$, dos vértices recém-alcançados (linha~\ref{hk:lk2}): o caminho segue obrigatoriamente pela aresta emparelhada, o que garante a alternância.
    \item \textbf{Busca dos caminhos disjuntos:} Primeiro, descartam-se os vértices da última camada que não são livres (linha~\ref{hk:del}), pois nenhum caminho aumentante termina neles. Em seguida, uma busca em profundidade parte de cada vértice livre $x_0 \in L_0$ (linha~\ref{hk:x0}) e avança sempre para um vizinho ainda não visitado da camada seguinte (linha~\ref{hk:adv}). Ao atingir a camada $k^*$, o caminho $P_k = (x_0, \dots, x_{k^*})$ é registrado (linha~\ref{hk:path}) e a busca a partir dessa raiz termina; se a busca fica sem vizinhos antes disso, recua uma camada (linha~\ref{hk:back}). Como cada vértice é visitado no máximo uma vez por fase, os caminhos encontrados são disjuntos em vértices, e o conjunto $\Delta M$ é maximal.
    \item \textbf{Complexidade:} Cada fase examina cada aresta um número constante de vezes, com custo $O(E)$. Após $\sqrt{V}$ fases, todo caminho aumentante restante tem mais de $\sqrt{V}$ arestas; como os caminhos aumentantes de $M \ominus M^*$ são disjuntos em vértices, restam menos de $\sqrt{V}$ deles, e cada fase seguinte elimina ao menos um. O número de fases é, portanto, $O(\sqrt{V})$, o que leva à complexidade total $O(E\sqrt{V})$ \cite{burkard}.
\end{itemize}

\subsubsection{Heurísticas}
\label{sec:h_emp_card_max_bip}

Por mais que o algoritmo de Hopcroft-Karp descrito em \ref{alg:hopcroft_karp} seja considerado relativamente eficiente, com complexidade $O(E \sqrt{V})$, ele ainda não é tão viável para instâncias muito grandes do problema de emparelhamento de cardinalidade máxima. Por isso, foram propostos alguns métodos aproximados (heurísticas) com o objetivo de resolver o problema com complexidade linear $O(E)$.

\paragraph{Heurística gulosa simples}
\label{p:h_gulosa_simples}

Um método que itera sobre as arestas do grafo em ordem arbitrária. Se a aresta $(u,v)$ conecta dois vértices livres, adicione essa aresta ao emparelhamento e remova os vértices $u$ e $v$. Tem complexidade de tempo linear $O(E)$ e garante um emparelhamento maximal, que conforme visto na seção \ref{p:emp_max_maximal}, é um emparelhamento ao qual não é possível adicionar mais arestas, mas que não é necessariamente o de maior cardinalidade. A solução também garante um emparelhamento de pelo menos $1/2$ (metade) da cardinalidade máxima \cite{vazirani2001}.

\begin{algorithm}[H]
    \DontPrintSemicolon
    \SetAlgoLined

    \caption{Heurística gulosa simples \cite{vazirani2001}}
    \label{alg:greedy_mcm}

    \KwIn{bipartite graph $G = (U, V; E)$}
    \KwOut{maximal matching $M$}
    $M := \emptyset$\;
    \For{each edge $(u, v) \in E$, in arbitrary order}{
        \If{$u$ and $v$ are both unmatched}{
            $M := M \cup \{(u, v)\}$\;
            mark $u$ and $v$ as matched\;
        }
    }
    \KwRet $M$\;
\end{algorithm}

\paragraph{Heurística de Karp-Sipser}

Heurística descrita por \citeonline{karp-sipser} para grafos esparsos que toma decisões com base na topologia do grafo ao invés de decisões aleatórias, tornando essas decisões mais seguras.

O algoritmo funciona reduzindo dinamicamente o grafo $G=(V,E)$ através da aplicação repetida de duas regras, nesta ordem de prioridade:

\begin{enumerate}
    \item \textbf{Redução de folhas de grau 1:} Se existir algum vértice $v$ no grafo atual com \textbf{grau 1} (apenas um vizinho $u$), a aresta $(u, v)$ é selecionada imediatamente para o emparelhamento. Isso é seguro pois, se $u$ é o único vizinho de $v$, ele também é a única opção para emparelhar $v$. Se isso não for feito, $v$ ficará isolado.
    \item \textbf{Escolha Aleatória:} Se \textbf{não} existirem vértices de grau 1 (ou seja, todos os vértices têm grau $\ge 2$), o algoritmo escolhe uma aresta $(u, v)$ aleatoriamente, adiciona ao emparelhamento e remove $u$ e $v$ do grafo. 
    
    Em grafos esparsos, isso tende a quebrar ciclos e formar novas folhas (vértices de grau 1), permitindo que o algoritmo volte a usar a regra 1. Porém, em grafos densos, essa escolha aleatória pode comprometer a precisão da solução.
\end{enumerate}

Para grafos esparsos, a regra 2 é raramente aplicada, pois a remoção de vértices na regra 1 frequentemente cria novas folhas, provocando um efeito cascata, fazendo com que a maior parte do grafo seja resolvida de maneira ótima. Com estruturas de dados adequadas, a complexidade do algoritmo é linear ($O(E)$). A implementação avaliada neste trabalho, por simplicidade, localiza os vértices de grau 1 por varredura, o que eleva sua complexidade a $O(V \cdot E)$.

\begin{algorithm}[H]
    \DontPrintSemicolon
    \SetAlgoLined

    \caption{Heurística de Karp-Sipser \cite{karp-sipser}}
    \label{alg:karp_sipser}

    \KwIn{bipartite graph $G = (U, V; E)$}
    \KwOut{maximal matching $M$}
    $M := \emptyset$\;
    \While{$E \neq \emptyset$}{
        \eIf{there is a vertex $v$ with $\deg(v) = 1$}{
            let $u$ be the only neighbor of $v$\;
            \AlgComment{safe: some maximum matching contains $(u, v)$}\;
        }{
            choose an arbitrary edge $(u, v) \in E$\;
            \AlgComment{the only step that may lose optimality}\;
        }
        $M := M \cup \{(u, v)\}$\;
        remove $u$, $v$ and all their incident edges from $G$\;
    }
    \KwRet $M$\;
\end{algorithm}

\begin{table}[htbp]
    \centering
    \caption{Algoritmos para emparelhamento de cardinalidade máxima em grafos bipartidos}
    \label{tab:emp_card_max_bip_algos}
    \renewcommand{\arraystretch}{1.2} % Adds a little vertical padding
    
    % Structure: l = left align, c = center, X = auto-wrap paragraph
    \begin{tabularx}{\textwidth}{@{} l l l l X @{}}
        \toprule
        \textbf{Algoritmo} & \textbf{Complexidade} & \textbf{Tipo} & \textbf{Avaliado} & \textbf{Descrição} \\
        \midrule
        
        Caminhos aumentantes & $O(VE)$ & Exato & Sim & 
        Busca iterativa (DFS/BFS) por caminhos alternantes entre vértices livres, invertendo arestas para aumentar a cardinalidade. \\
        \addlinespace 
        
        Edmonds-Karp & $O(VE^2)$ & Exato & Sim & 
        Implementação de Ford-Fulkerson que usa BFS para encontrar sempre o caminho aumentante mais curto (em nº de arestas) na rede residual. \\
        \addlinespace
        
        Dinic & $O(V^2E)$ & Exato & Não & 
        Algoritmo de fluxo que constrói grafos de níveis (BFS) e satura múltiplos caminhos simultaneamente via fluxo de bloqueio (DFS). Em redes de capacidade unitária, como a do emparelhamento bipartido, cai para $O(E\sqrt{V})$. \\
        \addlinespace
        
        Hopcroft-Karp & $O(E\sqrt{V})$ & Exato & Sim & 
        Busca múltiplos caminhos aumentantes disjuntos por fase usando BFS e DFS. \\
        \addlinespace
        
        Gulosa simples & $O(E)$ & Heurística & Sim & 
        Itera sobre arestas arbitrariamente, selecionando aquelas que conectam dois vértices livres. \\
        \addlinespace

        Karp-Sipser & $O(E)$ & Heurística & Sim* & 
        Prioriza remoção de vértices de grau 1; usa escolha aleatória quando não há folhas. \\
        
        \bottomrule
    \end{tabularx}
    \par\smallskip
    {\raggedright\footnotesize \textit{Avaliado}: incluído no ambiente de testes; a origem de cada implementação consta na Tabela~\ref{tab:proveniencia}. * A implementação avaliada localiza os vértices de grau 1 por varredura simples, com complexidade $O(V \cdot E)$.\par}
\end{table}

% TODO: possivelmente definir aplicações? exemplo: min-path-cover.md